%=============================================================================
\section{Results}
\label{sec:results}
%=============================================================================

\subsection{Cohort structure and group partitioning}
\label{subsec:dataset_chars}

Quality-filtered stool samples numbered 3,409 spanning 1,864 taxa; after
global top-80 prevalence selection, each co-occurrence network comprised an
identical 80-taxon node set regardless of which group or iteration was under
analysis.
Table~\ref{tab:groups} records group cardinalities before and after
demographic stratification matching.

\begin{table}[H]
\centering
\caption{%
  \textbf{Group sizes for each comparison (full and matched subsets).}
  Matching stratified on age-bin $\times$ sex $\times$ BMI-bin; samples
  missing any covariate were withheld from matched analyses.
  Bootstrap subsampling drew $n = 100$ per group per iteration.
}
\label{tab:groups}
\setlength{\tabcolsep}{8pt}
\begin{tabular}{llrrr}
\toprule
\textbf{Comparison} & \textbf{Group} & \textbf{Full $n$}
  & \textbf{Matched $n$} & \textbf{Subsample/iter} \\
\midrule
\multirow{2}{*}{IBD}
  & UC + Crohn's    & 165      & 135   & 100 \\
  & Healthy         & 3{,}084  & 404   & 100 \\
\addlinespace
\multirow{2}{*}{Diet}
  & Plant-based     & 367      & 308   & 100 \\
  & Omnivore        & 2{,}747  & 922   & 100 \\
\addlinespace
\multirow{2}{*}{Antibiotics}
  & Recent ($\leq$6\,mo.) & 590    & 479     & 100 \\
  & Absent ($>$1\,yr)     & 2{,}271 & 1{,}284 & 100 \\
\bottomrule
\end{tabular}
\end{table}

\subsection{IBD is associated with pan-topological collapse of co-occurrence loop structure}
\label{subsec:ibd_results}

The paired bootstrap distributions for the IBD comparison are shown in
Figure~\ref{fig:ibd_bootstrap}.
Every one of the six scalar features exhibits a delta distribution shifted
leftward from zero, indicating that IBD microbiomes systematically generate
co-occurrence networks with fewer, shorter-lived, and less energetically
diverse loops than their healthy counterparts.
No feature was directionally equivocal.

\begin{figure}[H]
  \centering
  \includegraphics[width=\textwidth]{../figures/agp_ibd_v2.png}
  \caption{%
    \textbf{Paired bootstrap delta distributions: IBD (UC + Crohn's)
    vs.\ Healthy, $N = 80$ taxa.}
    Each panel plots the distribution of the iteration-level difference
    $\delta = f^{(\text{IBD})} - f^{(\text{Healthy})}$ across 200
    paired iterations ($n = 100$ per group per iteration).
    Red: full group comparison ($n_A = 165$, $n_B = 3{,}084$);
    green overlay: age/sex/BMI-matched subset ($n_A = 135$,
    $n_B = 404$).
    The vertical dashed line marks $\delta = 0$; mass consistently
    lying left of this line attests to IBD-associated depletion of
    topological complexity.
    Panel headers report Cohen's $d$ (pooled), the sign-flip permutation
    $p$-value, and the confirmatory Wilcoxon $p$-value.
    All six features reach FDR significance at $p_{\text{FDR}} \leq 0.005$
    (Table~\ref{tab:ibd_results}).
  }
  \label{fig:ibd_bootstrap}
\end{figure}

Numerical results appear in Table~\ref{tab:ibd_results}.
The three features with the largest effect sizes---total persistence
($d = -2.28$), max Betti-1 ($d = -2.07$), and mean lifetime ($d = -1.65$)---
indicate that the healthy microbiome sustains roughly twice the cumulative
topological weight, twice the peak cycle density, and 42\% longer average
cycle lifetimes than the \ibd\ microbiome.
The permutation floor of $p = 0.002$ (the smallest achievable at
500 permutations; $\lfloor 1/501 \rfloor = 0.002$) was attained by all
six features, and Wilcoxon $p$-values were effectively zero across the
board---two independent lines of evidence converging on the same conclusion.

\begin{table}[H]
\centering
\caption{%
  \textbf{IBD vs.\ Healthy: full-group topological comparison, $N = 80$ taxa.}
  Means are averaged across 200 paired bootstrap iterations.
  Primary inference rests on sign-flip permutation $p$ (500 shuffles);
  $p_{\text{FDR}}$ is the Benjamini--Hochberg adjusted value across 18 tests.
  Wilcoxon $p$ is secondary and confirmatory only ($\star$ = FDR-significant).
}
\label{tab:ibd_results}
\small
\setlength{\tabcolsep}{5.5pt}
\begin{tabular}{lrrrrrr}
\toprule
\textbf{Feature} & \textbf{Mean IBD} & \textbf{Mean Healthy}
  & \textbf{Cohen's $d$} & \textbf{Perm $p$}
  & \boldmath$p_{\text{FDR}}$ & \textbf{Wilcoxon $p$} \\
\midrule
H1 count             & 16.87 & 23.66 & $-$1.34 & 0.002 & 0.004$^\star$ & $<$0.001 \\
H1 entropy           & 2.403 & 2.749 & $-$1.23 & 0.002 & 0.004$^\star$ & $<$0.001 \\
Total persistence    & 0.397 & 0.795 & $-$2.28 & 0.002 & 0.004$^\star$ & $<$0.001 \\
Mean lifetime        & 0.024 & 0.034 & $-$1.65 & 0.002 & 0.004$^\star$ & $<$0.001 \\
Max lifetime         & 0.080 & 0.119 & $-$1.45 & 0.002 & 0.004$^\star$ & $<$0.001 \\
Max Betti-1          & 3.44  & 7.09  & $-$2.07 & 0.002 & 0.004$^\star$ & $<$0.001 \\
\bottomrule
\end{tabular}
\end{table}

Covariate matching left the picture unaltered: all six features remained
FDR-significant in the matched subsample ($n_A = 135$, $n_B = 404$), and
effect sizes were numerically indistinguishable from the full-group estimates
(matched total persistence $d = -2.22$; matched max Betti-1 $d = -2.32$).

For geometric intuition, Figure~\ref{fig:ibd_tda} displays persistence
diagrams and Betti-1 curves from a representative single-pass comparison;
the sparser diagram and attenuated Betti-1 peak for the IBD group
corroborate the bootstrap evidence at the level of individual network topology.

\begin{figure}[H]
  \centering
  \includegraphics[width=0.92\textwidth]{../figures/agp_ibd_tda.png}
  \caption{%
    \textbf{Illustrative single-comparison TDA output: IBD vs.\ Healthy
    ($n = 500$ per group, $N = 80$ taxa).}
    \textit{Top left and centre:} Persistence diagrams.
    Connected-component bars (\Hzero, blue circles) and loop bars (\Hone,
    orange triangles) are plotted at their (birth, death) coordinates;
    proximity to the diagonal denotes short-lived, numerically fragile
    features.
    The IBD diagram is visibly sparser in $H_1$ features and shows a
    tighter cluster near the diagonal, reflecting the reduced loop
    persistence quantified in Table~\ref{tab:ibd_results}.
    \textit{Top right:} Betti-1 curves.
    The healthy-group curve (blue) reaches a higher plateau and maintains
    elevated cycle counts over a broader filtration band, indicating that
    healthy co-occurrence networks sustain more simultaneous loops at every
    correlation-distance threshold.
    \textit{Bottom row:} Persistence landscape $\lambda_1(t)$ (left),
    topological summary bar chart (centre), and reported test statistics
    (right).
    This panel is illustrative; statistical inference derives exclusively
    from the paired bootstrap analysis (Figure~\ref{fig:ibd_bootstrap},
    Table~\ref{tab:ibd_results}).
  }
  \label{fig:ibd_tda}
\end{figure}

\subsection{Plant-based diet: fewer but more tenacious co-occurrence cycles}
\label{subsec:diet_results}

The plant-based vs.\ omnivore comparison revealed a more nuanced topological
signature than the IBD contrast (Figure~\ref{fig:diet_bootstrap}).

\begin{figure}[H]
  \centering
  \includegraphics[width=\textwidth]{../figures/agp_diet_v2.png}
  \caption{%
    \textbf{Paired bootstrap delta distributions: Plant-based vs.\ Omnivore,
    $N = 80$ taxa.}
    Red = full groups ($n_A = 367$, $n_B = 2{,}747$);
    green = matched subset ($n_A = 308$, $n_B = 922$).
    H1 count shows a negative delta (plant-based eaters generate
    \emph{fewer} total loops), while mean lifetime, max lifetime, and
    max Betti-1 are positive (the loops they do form are more persistent
    and reach higher simultaneous densities).
    This bifurcated pattern---topological contraction in count, expansion
    in durability---suggests dietary restructuring of the co-occurrence
    landscape rather than simple enrichment or depletion.
  }
  \label{fig:diet_bootstrap}
\end{figure}

Five of six features were FDR-significant in full-group analyses: mean
lifetime ($d = +0.65$, $p_\text{FDR} = 0.004$), max Betti-1
($d = +0.46$, $p_\text{FDR} = 0.004$), total persistence
($d = +0.22$, $p_\text{FDR} = 0.030$), H1 count ($d = -0.34$,
$p_\text{FDR} = 0.007$), and H1 entropy ($d = -0.23$,
$p_\text{FDR} = 0.029$).
After matching, three features survived correction (mean lifetime,
max lifetime, max Betti-1), signalling partial---though not total---confounding
by demographic covariates.

\subsection{Antibiotic recency: an inconclusive topological signal}
\label{subsec:abx_results}

Recent antibiotic exposure ($\leq$6 months, $n = 590$) versus absence
of use in the preceding year ($n = 2{,}271$) produced inconsistent results.
In full-group analyses at $N = 80$, only max Betti-1 cleared FDR
($d = -0.39$, $p_\text{FDR} = 0.004$); in matched analyses at $N = 80$,
zero features survived correction.
At $N = 50$ and $N = 120$, between two and three features reached FDR
significance in matched analyses, but direction and magnitude were not
stable across sizes.
We decline to characterize this as a genuine topological effect.
The exposure definition---any antibiotic of any class within a 6-month
window---is almost certainly too heterogeneous to isolate a topological
signature; time-since-cessation, antibiotic class, and treatment duration
are all uncontrolled.

\subsection{Taxa-set sensitivity demonstrates IBD robustness and ABX fragility}
\label{subsec:sensitivity}

Figure~\ref{fig:sensitivity} displays Cohen's $d$ and FDR significance
across all three comparisons and all three taxon-set sizes for full groups.

\begin{figure}[H]
  \centering
  \includegraphics[width=\textwidth]{../figures/taxa_sensitivity_full.png}
  \caption{%
    \textbf{Taxa-set sensitivity heatmap: Cohen's $d$ by $N$ taxa (full groups).}
    Each cell encodes Cohen's $d$ for a feature (row) $\times$
    taxa-set size (column) combination; cells marked $\star$ cleared
    FDR correction (18-test BH, sign-flip permutation $p$-values).
    \textit{Centre panel (IBD):} All 18 cells are strongly negative
    and FDR-significant, confirming that neither the direction nor the
    significance of the \ibd\ finding depends on how many taxa are
    included ($d$ ranges from $-0.75$ to $-2.38$ across the 18 cells).
    \textit{Right panel (Diet):} Partial robustness; 3--4 of 6 features
    FDR-significant at each size, with consistent positive direction for
    lifetime metrics.
    \textit{Left panel (Antibiotics):} Erratic; effect-size direction
    varies across features and sizes, with no coherent pattern.
  }
  \label{fig:sensitivity}
\end{figure}

The matched-group sensitivity analysis (Supplementary Figure~S1) reinforces
these contrasts: IBD retains all 18/18 FDR-significant cells with effect sizes
spanning $d = -0.81$ to $d = -2.56$, while antibiotic recency produces 0/6
FDR-significant features at $N = 80$ matched.

Table~\ref{tab:sensitivity_summary} compresses these results into a
6-cell summary per subset.

\begin{table}[H]
\centering
\caption{%
  \textbf{Taxa-set sensitivity summary}: FDR-significant features out of
  6 per comparison and $N \in \{50, 80, 120\}$.
  Full groups (upper) and matched subsets (lower).
}
\label{tab:sensitivity_summary}
\small
\setlength{\tabcolsep}{10pt}
\begin{tabular}{llccc}
\toprule
\textbf{Subset} & \textbf{Comparison}
  & \boldmath$N = 50$ & \boldmath$N = 80$ & \boldmath$N = 120$ \\
\midrule
\multirow{3}{*}{Full}
  & IBD          & \textbf{6/6} & \textbf{6/6} & \textbf{6/6} \\
  & Diet         & 4/6 & 4/6 & 3/6 \\
  & Antibiotics  & 2/6 & 2/6 & 3/6 \\
\addlinespace
\multirow{3}{*}{Matched}
  & IBD          & \textbf{6/6} & \textbf{6/6} & \textbf{6/6} \\
  & Diet         & 4/6 & 3/6 & 4/6 \\
  & Antibiotics  & 3/6 & 0/6 & 2/6 \\
\bottomrule
\end{tabular}
\end{table}
